% sections/04_architecture.tex — Section IV (Architecture and Theoretical Grounding)
% Source: notes/paper_section_IV_draft.md v2 final, frozen 2026-05-12 evening
%   (Theorem 3' v3 reformulation + 4-fix reviewer audit integrated)
% Converted Day 5 2026-05-13.

\section{Architecture and Theoretical Grounding}
\label{sec:architecture}

\subsection{Implications of Theorem~\ref{thm:bound} for VLA Architecture Design}
\label{sec:architecture:implications}

Theorem~\ref{thm:bound} (\S\ref{sec:theory:bound}) provides an architecture-agnostic upper
bound on the language conditioning of any demonstration-supervised policy:
\begin{equation}
I(\Astar; \ell \mid o) \;\le\; I_{\dataD}(a; \ell \mid o) \;\le\; \eps.
\label{eq:bound_in_arch}
\end{equation}
By Remark~\ref{rem:arch_agnostic} of \S\ref{sec:theory:bound}, this bound applies uniformly
across the standard space of VLA action-decoder families: $L_2$ regression, $L_1$
regression, Conditional Flow Matching, Shortcut Flow Matching, and cross-entropy on
discretized actions. The bound is not bypassable through architectural choice within the
demonstration-supervised paradigm; bypassing it requires interventions outside the
demonstration loss itself (auxiliary language objectives, contrastive language alignment,
or inference-time conditioning on backbone-derived signals --- discussed in
\S\ref{sec:discussion}).

This observation reorients the question that architectural design can answer. The bound
determines what is \emph{achievable} in terms of policy-level language conditioning; the
architecture determines what failure mode manifests when the bound is tight. We distinguish
two architectural axes that operate independently of the bound.

\textbf{Decoder family.} The choice between deterministic action regression,
flow-matching-based generation, and cross-entropy over discretized actions does not change
the bound but does change the empirical character of policy behavior when language signal
is unavailable. Deterministic regression heads can still achieve high task SR by relying on
visual grounding alone (\S\ref{sec:results:oft_contrast}, OFT existence proof); they appear
to ``work'' at the SR level while satisfying the bound at the language-MI level.
Flow-matching-based decoders, in contrast, exhibit a distinctive sample-quality collapse
under language-orthogonal demonstrations: the matched-flow representation cannot learn to
discriminate among task-conditioned action distributions when the conditioning signal is
absent in the supervision data (\S\ref{sec:results:cfm_collapse} documents this empirically).

\textbf{Conditioning input richness.} Independent of decoder family, the set of inputs
supplied to the decoder --- observation features, proprioception, mesh-derived geometry,
grasp affordance maps, instruction embeddings --- determines what information the decoder
\emph{could} use to discriminate among tasks. Theorem~\ref{thm:bound} constrains the
conditioning that flows through the \emph{language channel}; it does not constrain
conditioning that flows through observation-derived channels. Component ablations
(\S\ref{sec:results:component_ablations}) systematically remove conditioning inputs to
probe the decoder's sensitivity along this axis.

These two axes motivate the architectural design of OTP-Soft, described next, and they
motivate the comparative evaluation strategy of \S\ref{sec:results}: OFT contrast as a
cross-architecture qualitative comparison (\S\ref{sec:results:oft_contrast}), and component
ablations for conditioning-input effects within the OTP-Soft architecture
(\S\ref{sec:results:component_ablations}).

\subsection{OTP-Soft Architecture}
\label{sec:architecture:otp_soft}

OTP-Soft inherits the OpenVLA-OFT pretrained backbone and replaces the action head with a
structured pipeline producing object-pose trajectories followed by language-agnostic action
decoding.

\paragraph{Backbone.}
We use the OpenVLA-OFT checkpoint~\cite{kim2025openvla_oft}, a 7.6B-parameter
vision-language model with a SigLIP+DINOv2 dual-encoder visual frontend fused to a
6-channel representation internally, and a LLaMA-2-derived language model trunk inherited
from OpenVLA~\cite{kim2024openvla}. We freeze the backbone (92M trainable parameters under
our training configuration are confined to downstream modules). Following the OFT
integration pattern, we access internal modules (\texttt{vision\_backbone},
\texttt{projector}, \texttt{llm\_backbone}) directly rather than through the outer forward
path, bypassing OFT's discretized action mask logic.

\paragraph{Object-pose trajectory head.}
A 138.3M-parameter trajectory head consumes backbone hidden states and emits an 8-step
trajectory of object poses ($N_{\text{obj}} \times H \times 7$ for $H = 8$ horizon steps
and 7-DoF pose per object). The head is trained with the shortcut model framework of
Frans et al.~\cite{frans2024shortcut} applied to the flow-matching
variant~\cite{lipman2023flow_matching} (denoted Shortcut Flow Matching throughout this
paper), chosen at design time for its accelerated convergence property under the shortcut
self-consistency objective relative to vanilla Conditional Flow Matching. We note that the
choice between Shortcut and vanilla CFM does not affect Theorem~\ref{thm:bound}
applicability (Remark~\ref{rem:arch_agnostic} of \S\ref{sec:theory:bound});
\S\ref{sec:results:cfm_collapse} documents that sample-quality collapse occurs under
Shortcut Flow Matching despite this design choice. We denote the head's stochastic output
as the latent trajectory $z$.

\paragraph{Language-agnostic decoder.}
The decoder $\phitheta$ accepts the trajectory $z$, proprioception, grasp affordance maps,
and object-geometry features, and emits a 7-DoF action chunk of horizon $H$. By
architectural construction, the instruction $\ell$ is excluded from the decoder's input
signature: $\ell$ does not appear in \texttt{forward()} arguments or as an input to any
submodule. This realizes Condition~\ref{cond:c3} of \S\ref{sec:theory:c3} at the
implementation level: $\ell \notin \mathrm{inputs}(\phitheta)$. The exclusion is enforced
by a \texttt{\_check\_forbidden} guard that raises \texttt{AssertionError} if any forbidden
name fragment appears among (i) \texttt{forward()} parameter names at module construction,
(ii) submodule and parameter names at construction, or (iii) \texttt{locals().keys()} and
\texttt{dir(self)} at every forward pass --- providing both static and runtime verification
of the C3 boundary.

The decoder also uses Shortcut Flow Matching as its loss family, matching the head. This
choice was made for consistency at design time and is empirically consequential:
\S\ref{sec:results:cfm_collapse} documents that the decoder inherits sample-quality
collapse independently of the head's behavior (oracle ablation in
\S\ref{sec:results:cfm_collapse} rules out the moving-target hypothesis).

\paragraph{Proprioception schema.}
Proprioception is encoded as an 8-dimensional vector: end-effector position (3) +
orientation quaternion (4) + a structurally zero gripper dimension (1). The zero gripper
dimension is by design --- gripper actions are predicted as part of the 7-DoF action
output rather than treated as a proprioception input, reflecting the LIBERO Panda
gripper's discrete open/close semantics. We frame this as an architectural choice rather
than a limitation: it preserves the action prediction's autonomy over gripper control
while keeping the proprioception schema dimensionally consistent across rollouts.

\paragraph{Mesh prior.}
Object geometry is encoded via a pose-invariant shape descriptor computed offline from
object meshes, forming a closed-set object library. Training and test splits in our
LIBERO-Spatial evaluation share five mesh identities. This is a deliberate scoping
choice --- pose-invariant geometry encoding allows the decoder to reason about object
affordances independently of pose without requiring per-test mesh generalization. The
implications of this closed-set assumption for cross-suite generalization are discussed in
\S\ref{sec:discussion}.

\subsection{Diagnostic Value of the OTP-Soft Architecture for This Study}
\label{sec:architecture:falsification}

The OTP-Soft architecture was originally designed as a candidate for addressing the
language conditioning concerns motivating this work. The theoretical analysis presented
here (\S\ref{sec:theory:bound}) shows that no architecture-internal pathway can reduce the
bound on policy-level language conditioning within the demonstration-supervised paradigm
--- a finding that is independent of any specific architecture. We retain OTP-Soft as the
experimental platform for this paper on the basis of its diagnostic properties rather
than its original architectural hypothesis:

\textbf{1.~Architectural separation enables component ablation.} The explicit C3 boundary
at the decoder, together with the well-defined conditioning input set $\{z,
\text{proprioception}, \text{grasp affordance}, \text{object geometry}\}$, allows
\S\ref{sec:results:component_ablations} to systematically drop conditioning channels and
isolate their contribution to decoder behavior. Architectures without such clean separation
(e.g., end-to-end transformer decoders consuming arbitrary token mixtures) would conflate
input-ablation effects with attention-allocation effects.

\textbf{2.~Decoder-family qualitative comparison with OFT.} OTP-Soft and OFT share the
same pretrained backbone but differ in multiple architectural dimensions: decoder family
(Shortcut Flow Matching vs. deterministic regression), intermediate representation (latent
trajectory $z$ vs. none), object query mechanism (5 object queries vs. none), conditioning
input set (mesh and grasp affordance vs. single-image input), and proprioception schema
(8-dim with structural zero vs. 7-dim). We cannot fully isolate the decoder-family
contribution from these other architectural differences within the present work. The
contrast does inform our qualitative claim in \S\ref{sec:results:oft_contrast} that at
least one demonstration-supervised architecture (OFT) achieves high SR on LIBERO-Spatial
while at least one other (OTP-Soft V3) fails on the same supervision protocol and same
backbone --- a contrast that complements but does not isolate the decoder-family axis.
Controlled architecture experiments designed to isolate this axis are left to future work.

\textbf{3.~Diagnostic infrastructure generalizes.} The Gate~1, M3, and C1 measurements
developed for OTP-Soft (\S\ref{sec:results:diagnostic}) do not depend on architectural
specifics beyond access to backbone hidden states, demonstration trajectories, and policy
latent representations. These measurements apply equally to OFT, RT-2, $\pi_0$, or other
VLA architectures sharing a comparable structural decomposition. The diagnostic
methodology is the principal contribution of this paper; OTP-Soft is its first
instantiation.

A reviewer might reasonably ask whether the architectural complexity of OTP-Soft is
justified given that the architecture does not solve the conditioning problem it was
originally designed to address. Our answer is twofold. First, the architecture's
diagnostic value (above three reasons) holds independently of its original design
hypothesis --- it serves as a controlled experimental platform regardless of original
motivation. Second, the absence of an architectural fix in this work --- established
theoretically by Theorem~\ref{thm:bound} and instantiated empirically by the
\S\ref{sec:results:diagnostic} Mantel measurements and \S\ref{sec:results:cfm_collapse}
sample-quality test --- is itself a substantive finding. Future work should pursue
auxiliary-objective interventions outside the demonstration-supervised loss
(\S\ref{sec:discussion}), informed by the bound this paper establishes.
